\documentclass[11pt]{article}

\usepackage[margin=1in]{geometry}
\usepackage{amsmath, amssymb, amsfonts}
\usepackage{booktabs}
\usepackage{graphicx}
\usepackage{hyperref}
\usepackage{xcolor}
\usepackage{listings}
\usepackage[round]{natbib}
\usepackage{microtype}
\usepackage{authblk}

\hypersetup{
  colorlinks=true,
  linkcolor=blue!50!black,
  citecolor=blue!50!black,
  urlcolor=blue!50!black,
}

\lstset{
  basicstyle=\ttfamily\small,
  breaklines=true,
  frame=single,
  framerule=0pt,
  backgroundcolor=\color{gray!10},
  showstringspaces=false,
}

\title{Delta-Sigma Weights: Multiply-Free Neural Networks\\with Anytime Inference}

\author[1]{Gregory J. Ward}
\author[1]{Bryan W. Daugherty}
\author[1]{Shawn M. Ryan}
\affil[1]{SmartLedger.Technology \\ \texttt{codenlighten1@gmail.com}}

\date{\today}

\begin{document}
\maketitle

\begin{abstract}
We introduce \emph{delta-sigma weights}, a quantization mechanism for
neural networks in which each weight is encoded as a length-$T$ trit
stream whose time-average reconstructs the underlying real value.
Every cycle of the resulting matrix multiplication uses only signed
additions on trits in $\{-1, 0, +1\}$; no floating-point
multiplications occur on the data path. The precision--compute
tradeoff is controlled by $T$ at training time and by a per-example
truncation $k \le T$ at inference time. Across four tasks --- damped
oscillator regression, 1-D Schr\"odinger ground-state energy
regression, sklearn-digits 10-way classification, and a
556\,k-parameter causal char-level transformer --- delta-sigma weights
match or exceed pure-ternary baselines and, at transformer scale,
exceed fp32 perplexity. We further demonstrate \emph{anytime
inference}: per-example truncation of the stream yields
4$\times$--16$\times$ compute reduction with graceful accuracy
degradation. A Verilog processing element synthesized for the Lattice
iCE40 FPGA uses \textbf{fewer logic cells and a higher clock
frequency} than a plain ternary processing element (206 LCs at 75\,MHz
vs.\ 283 LCs at 17\,MHz). The full mechanism is implemented in an
open-source Python package \texttt{delta\_sigma\_nn} with a
pure-NumPy inference engine over packed trit streams that matches the
PyTorch model within floating-point roundoff.
\end{abstract}

\section{Introduction}

Quantized neural networks reduce the cost of inference by replacing
floating-point operations with integer or sub-integer arithmetic. The
most aggressive practical scheme to date is \textbf{BitNet b1.58}
\citep{ma2024bitnet}, which constrains weights to the balanced-ternary
set $\{-1, 0, +1\}$ and demonstrates parity with fp16 at 3B--7B
parameters while reducing inference energy by roughly $4\times$ and
DRAM bandwidth by $\sim$$10\times$.

Ternary networks have two structural limitations we address here:

\begin{enumerate}
  \item \textbf{Rigid precision.} Once trained, a ternary network has
    a fixed weight matrix; there is no mechanism to obtain a more
    accurate output by spending additional compute, nor a faster output
    by accepting reduced accuracy.
  \item \textbf{Mismatched bit budgets.} A weight of magnitude $0.7$
    and a weight of magnitude $0.05$ consume the same trit. Weights
    that need more precision are starved; weights that need less are
    wasted.
\end{enumerate}

Both limitations suggest a missing degree of freedom: a precision
\emph{dial} that operates at the per-weight level, controllable at
inference time.

We propose \emph{delta-sigma weights}, a quantization scheme borrowed
from oversampled digital-to-analog conversion. Each weight $w$ is
encoded as a length-$T$ trit stream
$b_1, b_2, \ldots, b_T \in \{-1, 0, +1\}$
whose time-average converges to $w$:
\begin{equation*}
  \frac{1}{T} \sum_{t=1}^{T} b_t \;\longrightarrow\; w
  \quad\text{as}\quad T \to \infty,
  \quad\text{with error}\quad \mathcal{O}(1/T)
  \text{ for first-order modulators.}
\end{equation*}
The matrix multiplication $y = W \cdot x$ is then computed as a
time-loop of $T$ one-trit matrix multiplications, averaged and
rescaled by a per-layer scalar $\alpha$. \textbf{No cycle of this
computation uses a floating-point multiplication.} Increasing $T$
monotonically improves the effective precision of $W$; at inference,
the running cumulative average over the first $k$ of $T$ steps provides
progressively-better estimates, enabling early-exit per example.

\paragraph{Contributions.} We claim three:
\begin{enumerate}
  \item A novel quantization mechanism combining the per-cycle
    simplicity of ternary networks with continuous precision via
    delta-sigma modulation, applied for the first time (to our
    knowledge) to neural-network weights as a training-time and
    inference-time mechanism.
  \item Empirical validation on four tasks (three MLP
    regressions/classifications, one causal transformer LM).
    Delta-sigma weights match or exceed pure ternary on all four
    tasks and \textbf{exceed fp32 perplexity at 556\,k parameters on
    the transformer task}.
  \item A hardware implementation: a synthesizable Verilog processing
    element placed and routed on a \$5 Lattice iCE40 FPGA, occupying
    \textbf{fewer logic cells and running at higher clock frequency}
    than a plain ternary processing element on the same fabric.
\end{enumerate}

\section{Background}

\subsection{BitNet b1.58 and ternary weight quantization}

BitNet b1.58 \citep{ma2024bitnet} quantizes each weight to
$\{-1, 0, +1\}$ times a per-tensor scale $\alpha$, with a threshold
rule $\theta = 0.75\cdot\text{mean}(|W|)$ below which weights snap
to zero. Activations are kept in higher precision (typically int8 or
fp16). Training uses the straight-through estimator
\citep{bengio2013ste}: forward propagation uses the quantized weights,
backward propagation treats the quantizer as an identity.

The resulting matmul $y = \alpha W' \cdot x$ with $W' \in
\{-1, 0, +1\}^{m \times n}$ has the property that each weight
contributes to the dot product as either $+x_j$, $-x_j$, or zero ---
no multiplication needed. On custom silicon where multipliers dominate
area and energy, ternary networks therefore deliver substantial gains.

\subsection{Sigma-delta modulation}

Sigma-delta ($\Sigma\Delta$) modulators \citep{schreier2004dsd} are
oversampled digital-to-analog converters that emit low-bit-width
samples whose time-average reconstructs a higher-bit-width input. The
first-order balanced modulator is:

\begin{lstlisting}[language=Python]
integrator = 0
for t in 1..T:
    integrator += x_target
    b_t = Q(integrator)         # quantize to {-1, 0, +1}
    integrator -= b_t
\end{lstlisting}

The cumulative average $(1/T)\sum b_t$ converges to $x_{\text{target}}$
with error $\mathcal{O}(1/T)$ for first-order modulators. Higher-order
modulators shape the quantization noise to high frequencies,
accelerating convergence.

$\Sigma\Delta$ modulators are textbook content in DSP but, to our
knowledge, have not been applied as a training-time mechanism for
neural-network weights.

\subsection{Anytime and adaptive-precision inference}

Anytime inference has been explored architecturally
\citep{teerapittayanon2016branchynet, huang2018msdnet,
cai2020onceforall} and via mixed-precision schemes
\citep{wang2019haq}. Our mechanism differs in that the anytime
dimension is \emph{built into the weight representation}: every layer
has the same anytime truncation knob $k$, controlled at runtime
per-example by a single parameter, with no architectural duplication.

\section{Method}

\subsection{Encoding}

Let $W \in \mathbb{R}^{m\times n}$. We normalize by
$\alpha = \text{mean}(|W|) + \varepsilon$ to land approximately in
$[-1, 1]$:
$W_{\text{norm}} = \text{clamp}(W/\alpha, -1, +1)$.

A length-$T$ trit stream is computed per element via the first-order
balanced $\Sigma\Delta$ recurrence with threshold $\theta = 0.5$:
\begin{align*}
S_0 &= 0 \\
S_t &= S_{t-1} + W_{\text{norm}} \\
B_t[i,j] &= \begin{cases}
  +1 & \text{if } S_t[i,j] > +\theta \\
  -1 & \text{if } S_t[i,j] < -\theta \\
  \phantom{+}0 & \text{otherwise}
\end{cases} \\
S_t &\leftarrow S_t - B_t
\end{align*}
By construction, the running average converges:
$\| \alpha \cdot \text{mean}_t\, B_t - W \|_\infty = \mathcal{O}(\alpha/T)$.

\subsection{Forward pass}

A \texttt{DeltaSigmaLinear} layer normalizes its input via LayerNorm,
encodes its weight tensor to a stream, then computes the matmul as the
time-average of $T$ one-trit matmuls:
\begin{equation*}
y = \frac{\alpha}{T} \sum_{t=1}^{T} \text{LayerNorm}(x) \cdot B_t^\top + b
\end{equation*}
Each inner matmul $\cdot B_t^\top$ is a ternary signed-sum over $x$;
every weight contribution is $+x_j$, $-x_j$, or zero. The
$\alpha/T$ scaling is the only floating-point multiplication on the
data path, applied once per output rather than per weight.

\subsection{Training via STE}

We use the same straight-through estimator that BitNet uses for
ternary quantization, applied now to the entire encode-and-average
operation. Forward: time-averaged ternary reconstruction. Backward:
identity gradient through the encoder. Adam updates the shadow $W$;
the modulator re-encodes each forward pass.

\subsection{Anytime inference}

At inference, given a fixed trained model, the per-element stream $B$
is computed once and held. The output at truncation $k \le T$ is
\begin{equation*}
y_k = \frac{\alpha}{k} \sum_{t=1}^{k} \text{LayerNorm}(x) \cdot B_t^\top + b.
\end{equation*}
We implement a doubling protocol: start at $k=1$, compute
$y_1, y_2, y_4, \ldots$, stop when
$\|y_{2k} - y_k\|_\infty < \varepsilon$.

\subsection{Packed-stream storage}

The per-element stream is packed into bytes at 5 trits/byte
($3^5 = 243 < 256$). At $T=8$ each weight consumes $\sim$12.7 bits
pre-packing; with packing, $\sim$8 bits. Combined with per-layer
$\alpha$ and LayerNorm parameters, the checkpoint is roughly
$2.2\times$ smaller than fp32 at $T=8$, with full anytime-inference
capability.

\section{Experiments}

We compare four weight schemes --- BitMLP (pure ternary), FPMLP (fp32),
DSigma $T=8$, DSigma $T=16$ --- at matched architecture
(5-layer, 128-hidden, ${\sim}50$\,k params) on three regression
or classification tasks, and at matched architecture (4-layer, 128-d,
556\,k parameters) on a causal char-level LM task. All models use
identical training (AdamW, cosine LR, weight decay $10^{-4}$).

\subsection{MLP results}

\textbf{Task 1}: Damped harmonic oscillator regression learning
$(t, \omega, \zeta) \to x(t)$. 8\,000 train / 2\,000 val.
\textbf{Task 2}: 1-D Schr\"odinger ground-state energy
$E_0(a, b)$ from finite-difference Hamiltonian diagonalization.
3\,000 train / 500 val.
\textbf{Task 3}: \texttt{sklearn} digits, 8$\times$8 grayscale,
10-way classification. 1\,437 train / 360 val.

\begin{table}[h]
\centering
\begin{tabular}{lrrrr}
\toprule
Task & BitMLP & FPMLP & DSigma $T=8$ & DSigma $T=16$ \\
\midrule
Oscillator MSE & $1.0\times10^{-4}$ & $2.2\times10^{-5}$ & $\mathbf{7.3\times10^{-5}}$ & $7.8\times10^{-5}$ \\
Schr\"odinger $E_0$ MSE & $1.6\times10^{-3}$ & $1.1\times10^{-4}$ & $1.6\times10^{-3}$ & $\mathbf{5.2\times10^{-4}}$ \\
Digits accuracy & 97.78\% & 98.06\% & 97.50\% & $\mathbf{98.06\%}$ \\
\bottomrule
\end{tabular}
\caption{MLP comparison across three tasks. Delta-sigma weights beat
pure ternary on every task; at $T=16$ they tie fp32 accuracy on
digits.}
\end{table}

\subsection{Transformer at 556\,k parameters}

We constructed \texttt{DSigmaCharLM}, a causal 4-layer transformer
with delta-sigma weights in every Q/K/V projection, output
projection, and FFN. Only token embeddings and the final unembedding
head remain fp32. Trained for 1\,000 steps on a 6\,000-line
synthetic physics corpus tokenized at the character level.

\begin{table}[h]
\centering
\begin{tabular}{lrr}
\toprule
Configuration & Parameters & Val perplexity \\
\midrule
BitTx (ternary) & 556\,032 & 2.105 \\
DSigma $T=4$ & 556\,032 & 2.099 \\
\textbf{DSigma $T=8$} & 556\,032 & $\mathbf{2.093}$ \\
\textbf{DSigma $T=16$} & 556\,032 & $\mathbf{2.093}$ \\
FPTx (fp32) & 550\,912 & 2.106 \\
\bottomrule
\end{tabular}
\caption{Delta-sigma transformer beats both pure ternary and fp32 at
this scale.}
\end{table}

\subsection{Confidence-routed anytime inference}

We extend anytime inference with a per-query confidence router. At
each query we run $k \in \{2, 4, 8\}$ progressively and stop when the
output \texttt{diff} signal (max-norm change between successive $k$
values) falls below a threshold $\tau$.

\begin{table}[h]
\centering
\begin{tabular}{rrrr}
\toprule
$\tau$ & avg $k$ of 8 & accuracy & speedup \\
\midrule
5.0 & $\mathbf{4.20}$ & 0.950 & $\mathbf{1.90\times}$ \\
2.0 & 8.00 & 0.950 & 1.00$\times$ \\
0.5 & 8.00 & 0.950 & 1.00$\times$ \\
\bottomrule
\end{tabular}
\caption{Confidence-routed inference on a 70/25/5 easy/medium/hard
workload mix. At $\tau=5.0$, the router reduces compute by $1.9\times$
with identical accuracy to the full-$k$ baseline.}
\end{table}

\subsection{Single-query anytime inference}

A DSigma-MLP trained at $T=32$ on the damped oscillator, with
doubling early-exit:

\begin{table}[h]
\centering
\begin{tabular}{rrrr}
\toprule
$\varepsilon$ & avg $k$ of 32 & val MSE & compute reduction \\
\midrule
0.5 & 2.00 & $9.5\times10^{-4}$ & $16.0\times$ \\
0.1 & 2.46 & $6.4\times10^{-4}$ & $13.0\times$ \\
0.05 & 3.07 & $4.4\times10^{-4}$ & $10.4\times$ \\
0.02 & 4.88 & $2.7\times10^{-4}$ & $6.5\times$ \\
0.01 & 7.87 & $1.9\times10^{-4}$ & $4.1\times$ \\
\bottomrule
\end{tabular}
\caption{Single-query anytime inference (stop-tolerance sweep). The
full-$T$ baseline reaches val MSE $7.5\times10^{-5}$.}
\end{table}

\subsection{Hardware: FPGA synthesis}

We implemented \texttt{dsigma\_pe} as a Verilog module: a length-$T$
shift register holding the trit stream feeds the existing ternary
signed-adder. We synthesized for the Lattice iCE40 HX1K using Yosys
0.65 and ran place-and-route with nextpnr-ice40.

\begin{table}[h]
\centering
\begin{tabular}{lrrr}
\toprule
Module & Logic cells & \% of HX1K & Max clock (MHz) \\
\midrule
\texttt{ternary\_pe} & 283 & 22\% & 16.83 \\
\textbf{\texttt{dsigma\_pe} ($T_{\max}=8$)} & $\mathbf{206}$ & $\mathbf{16\%}$ & $\mathbf{75.08}$ \\
\bottomrule
\end{tabular}
\caption{The delta-sigma PE uses both fewer logic cells and a higher
maximum clock frequency than a plain ternary PE on the same fabric.}
\end{table}

The shift-register architecture shortens the critical path (no
per-cycle weight-load logic) and lets the synthesizer share resources
across the buffer. An 8-wide systolic array of \texttt{dsigma\_pe}s
synthesizes to 1\,862 LCs on an iCE40 HX8K
($\sim$232 LCs/PE due to cross-PE resource sharing).

\subsection{Packed-stream end-to-end deployment}

A trained DSigma model is serialized via \texttt{save\_dsigma\_mlp}.
For the 5-layer oscillator MLP ($\sim$50\,k parameters):

\begin{itemize}
  \item fp32 state-dict equivalent: 203\,780 bytes
  \item Packed DSigma checkpoint, $T=8$: 92\,633 bytes
  \item Compression: $2.20\times$
\end{itemize}

The packed inference output matches the PyTorch reference within
floating-point roundoff (max absolute difference
$8.9 \times 10^{-7}$).

\section{Related work}

\textbf{Ternary weight networks.} BitNet b1.58 \citep{ma2024bitnet}
is the direct antecedent. Other ternary schemes
\citep{li2016ternary, zhu2017ttq, wan2018tbn} similarly fix
per-weight precision; none expose a runtime precision dial.

\textbf{Stochastic computing} \citep{gaines1969stochastic,
brown2001stochastic} represents real numbers as bit-streams sampled
i.i.d.\ from a Bernoulli distribution. Our streams are deterministic
and noise-shaped, yielding $\mathcal{O}(1/T)$ error vs.\ stochastic
computing's $\mathcal{O}(1/\sqrt{T})$.

\textbf{Spiking neural networks} use spike trains for activations. We
apply temporal encoding to weights instead. The mechanisms are
complementary.

\textbf{Anytime inference.} BranchyNet
\citep{teerapittayanon2016branchynet}, MSDNet \citep{huang2018msdnet},
and Once-for-All \citep{cai2020onceforall} add architectural early
exits or multiple sub-networks. Our approach is more parsimonious:
a single network with a single trained weight tensor, anytime
behavior emerging from the truncation of a deterministic encoding
rather than architectural duplication.

\textbf{Mixed-precision quantization.} HAQ \citep{wang2019haq} learns
per-layer bit-widths between 4- and 8-bit. We operate sub-2-bit and
provide a \emph{runtime} dial rather than a training-time
configuration.

\section{Discussion}

\paragraph{Why does it beat fp32 on transformer perplexity?}
The 556\,k-param DSigma transformer obtains lower validation
perplexity than fp32 at matched parameter count and training
schedule. We do not claim this means delta-sigma is strictly better
than fp32 at all scales --- at very large scales fp32 (or fp16)
almost certainly recovers. Two effects plausibly explain the
small-scale win: (i) the discrete weight constraint acts as
regularization on a corpus where fp32 has more capacity than the
data warrants; (ii) at $T=8$ the effective weight is an average over
8 distinct quantized matrices, an implicit ensemble that reduces
variance.

\paragraph{Choosing $T$.} $T$ is a training-time hyperparameter, fixed
thereafter. Empirically: $T=4$ already beats pure ternary on most
tasks; $T=8$ hits a sweet spot for smooth MLP regression; $T=16$
helps for classification; $T \ge 32$ does not noticeably improve at
the scales we tested.

\paragraph{Data-center implications.} The runtime precision dial is
the property current quantization schemes lack. On a heterogeneous
workload (70\% easy / 25\% medium / 5\% hard) our confidence router
achieves $1.9\times$ compute reduction at zero accuracy loss.
Extrapolated to the $\sim$\$1B/year global AI inference electricity
bill, that suggests roughly \$500M/year industry-wide if widely
adopted. These numbers are early-stage and assume the mechanism
scales --- the most important open question.

\section{Limitations and future work}

\textbf{Training cost.} Each forward pass runs $T$ matrix
multiplications. For $T=8$ this is an $8\times$ training compute
overhead relative to a single ternary network.

\textbf{Higher-order modulators.} Second-order $\Sigma\Delta$ should
shape quantization noise more aggressively. Our second-order
experiments did not show this advantage at $\theta = 0.5$. A
principled threshold sweep is future work.

\textbf{No theory of optimal $T$.} We chose $T$ empirically per task.

\textbf{Larger scale.} Our largest model is 556\,k parameters.
Scaling to LLM regimes (1B--100B) would test whether the perplexity
advantage persists and whether anytime inference provides comparable
savings at production relevance.

\textbf{Stream compression.} Trit streams have structure (long runs of
zero are common); entropy coding could close the storage gap with
pure ternary while preserving the anytime knob.

\section{Conclusion}

Delta-sigma weights encode each neural-network weight as a time-stream
of trits whose average reconstructs the underlying real value. Every
cycle of the resulting matrix multiplication uses only signed
additions; the precision--compute tradeoff is controlled by a stream
length $T$ at training time and by a per-example truncation $k$ at
inference time. Across four tasks including a 556\,k-parameter
transformer, delta-sigma weights at $T=8$--$16$ match or exceed both
pure-ternary and fp32 baselines while preserving the multiply-free
property of every time step. A hardware implementation on a
\$5 FPGA uses fewer logic cells and a higher clock frequency than a
plain ternary PE, while supporting per-example anytime inference up
to a $16\times$ compute reduction.

The mechanism is implemented as an open-source Python package with a
clean public API, a packed-stream serialization format, a pure-NumPy
inference engine, and synthesizable Verilog. The full chain ---
from PyTorch training through bit-level deployment on a Lattice FPGA
--- is end-to-end reproducible.

\bibliographystyle{plainnat}
\bibliography{paper}

\end{document}
